\section*{अध्याय \ref{ch:b2:ore-deposits} --- अयस्क निक्षेप}

\begin{solution}{exo:b2:ore-deposits:1}
(क) $\qty{0.5}{\percent} = \qty{5000}{ppm}$, और $5000/28 = 179$। (ख) $\qty{6}{g/t} = \qty{6000}{ng/g}$,
और $6000/1.5 = 4000$। (ग) क्रोमाइट में द्रव्यमान से $2 \times 52.00/(55.85 + 104.00 + 64.00) = 0.4646$
क्रोमियम है; क्रोमाइटाइट में $0.9 \times 0.4646 = 0.418$, अर्थात् \qty{418000}{ppm}, और
$418000/92 = 4540$।
\end{solution}

\begin{solution}{exo:b2:ore-deposits:2}
(क) $914 \times 0.0053 = \qty{4.84}{Mt}$ ताँबा और $914 \times 10^6 \times 0.40 = \num{3.66e8}$ g,
अर्थात् \qty{366}{t} सोना; $1640 \times 0.0044 = \qty{7.22}{Mt}$ और \qty{738}{t}। (ख) $1640 - 914 =
\qty{726}{Mt}$, जिनमें $7.216 - 4.844 = \qty{2.37}{Mt}$ ताँबा है, अर्थात् माध्य ग्रेड
$2.372/726 = \qty{0.327}{\percent}$: अतिरिक्त खंड \qty{0.3}{\percent} कट-ऑफ़ से नीचे
और उसके ठीक ऊपर हैं। (ग) वे दो कट-ऑफ़ पर प्रश्न का उत्तर देते हैं; कट-ऑफ़ के बिना दोनों में से कोई भी
निक्षेप का आकार नहीं है।
\end{solution}

\begin{solution}{exo:b2:ore-deposits:3}
(क) $(19.3 - 1.00)/(2.65 - 1.00) = 18.3/1.65 = 11.1$। (ख) समान गतियों के लिए $(\rho_g - 1)d_g^2 =
(\rho_q - 1)d_q^2$ चाहिए, अतः $d_g/d_q = \sqrt{1.65/18.3} = 0.300$: सोने का कण अपने से लगभग तीन गुना
बड़े क्वार्ट्ज़ कण की तरह बैठता है। (ग) $d_q = 0.1\sqrt{3.7/1.65} = \qty{0.150}{mm}$।
\end{solution}

\begin{solution}{exo:b2:ore-deposits:4}
प्रति इकाई क्षेत्रफल और घनत्व, परत अपना ताँबा रखती है और निक्षालित स्तंभ का ताँबा पाती है:
$g = g_0(h_e + \eta h_l)/h_e$। (क) $0.4 \times (50 + 150)/50 = \qty{1.6}{\percent}$।
(ख) $0.4 \times (50 + 0.8 \times 150)/50 = \qty{1.36}{\percent}$। (ग) $h_l = 50 \times (1.7/0.4 - 1)
= \qty{163}{m}$। ये ग्रेड संवर्धित पोर्फ़िरी परतों के 0.4--\qty{1.7}{\percent} परास में
आते हैं।
\end{solution}

\begin{solution}{exo:b2:ore-deposits:5}
(क) प्लैटिनम: $C_S = 10 \times 10^4 (R + 1)/(R + 10^4)$ \unit{ppb}: \qty{1.00}{ppm} ($R = 100$),
\qty{9.10}{ppm} (\num{1000}), \qty{50.0}{ppm} (\num{10000}), \qty{90.9}{ppm} (\num{100000})।
निकेल: $C_S = 200 \times 500 (R + 1)/(R + 500)$ \unit{ppm}: \qty{1.68}{\percent}, \qty{6.67}{\percent},
\qty{9.52}{\percent}, \qty{9.95}{\percent}। (ख) $D\,C_L$: \qty{100}{ppm} प्लैटिनम,
\qty{10}{\percent} निकेल। (ग) निकेल अपनी अधिकतम सीमा के \qty{5}{\percent} के भीतर $R = \num{10000}$ से आगे
रहता है, इसलिए $R$ के दस गुना बदलने पर भी उसकी मात्रा बहुत कम बदलती है; वहाँ प्लैटिनम अब भी लगभग दोगुना हो रहा
होता है, इसलिए उसकी मात्रा $R$ को मापती है।
\end{solution}

\begin{solution}{exo:b2:ore-deposits:6}
(क) $\mu_M = 2.373 - 0.364 - 2 = 0.009$; $\sigma_M = 0.702$; माध्यिका $10^{0.009} = \qty{1.02}{Mt}$।
(ख) $z = (1 - 0.009)/0.702 = 1.412$, $1 - \Phi(1.412) = 0.0790$: 422 में से लगभग 33 निक्षेप।
(ग) $\sigma = \sqrt{0.654^2 + 0.199^2} = 0.684$, $z = 1.449$, $1 - \Phi = 0.0737$, लगभग 31: दुर्बल
सहसंबंध कुछ विशाल निक्षेप जोड़ देता है।
\end{solution}

\begin{solution}{exo:b2:ore-deposits:7}
सभी खंडों का माध्य $0.40\,e^{0.18} = \qty{0.479}{\percent}$। $c = 0.3$ के लिए: $z = \ln(0.40/0.3)/0.60 =
0.4795$, $\Phi(z) = 0.6842$, $\Phi(z + s) = 0.8598$: $T = \qty{684}{Mt}$, $\bar g = 0.479 \times
0.8598/0.6842 = \qty{0.602}{\percent}$, $M = 1000 \times 0.00479 \times 0.8598 = \qty{4.12}{Mt}$।
$c = 0.6$ के लिए: $z = -0.6758$, $\Phi = 0.2496$, $\Phi(z + s) = 0.4698$: \qty{250}{Mt},
\qty{0.901}{\percent} पर, \qty{2.25}{Mt}। चित्र पर, \qty{0.3}{\percent} के कट-ऑफ़ पर नीला वक्र लगभग
\qty{680}{Mt} और लाल वक्र लगभग \qty{0.6}{\percent} पढ़ता है।
\end{solution}

\begin{solution}{exo:b2:ore-deposits:8}
(क) तरल: $\num{5e6}/(\num{1e-3} \times 0.5) = \qty{1.00e10}{t}$; \omterm{def:b2:magma-emplacement:magma}{मैग्मा}: $\num{1e10}/0.035 =
\qty{2.86e11}{t}$। (ख) $\num{2.86e14}/2600 = \qty{1.10e11}{m^3}$, \qty{110}{km^3}, लगभग
\qty{4.8}{km} भुजा वाला \omterm{def:b2:magma-emplacement:magma}{मैग्मा} का एक घन। (ग) \omterm{def:b2:magma-emplacement:magma}{मैग्मा} में $\num{2.86e11}
\times \num{50e-6} = \qty{1.43e7}{t}$ ताँबा था, और अयस्क में उसका $5/14.3 = 0.350$ भाग।
\end{solution}

\begin{solution}{exo:b2:ore-deposits:9}
(क) लघुगणकों का माध्य 1.411, मानक विचलन 0.0664; माध्यिका $10^{1.411} = \qty{25.8}{ppm}$;
देहली $10^{1.411 + 0.133} = 10^{1.544} = \qty{35.0}{ppm}$। (ख) $(4 \times 400 + 36 \times 25.8)/40 =
\qty{63.2}{ppm}$, जो \qty{35.0}{ppm} से ऊपर है: स्वच्छ अवसाद के नौ भागों से तनुकरण के बाद भी
विसंगत। (ग) $4 \times 400 + (A - 4) \times 25.8 = 35.0\,A$ से $A = 4 \times (400 -
25.8)/(35.0 - 25.8) = \qty{163}{km^2}$ मिलता है।
\end{solution}

\begin{solution}{exo:b2:ore-deposits:10}
(क) $e^{su}\,e^{-u^2/2} = e^{-(u - s)^2/2 + s^2/2}$; $\sqrt{2\pi}$ से भाग दीजिए। (ख) $E[e^{\mu + su}] =
e^{\mu}\int e^{su}\varphi(u)\,\mathrm{d}u = e^{\mu + s^2/2}\int\varphi(u - s)\,\mathrm{d}u = e^{\mu +
s^2/2}$। (ग) $\Phi$ वर्धमान है और $s > 0$, अतः $\Phi(z + s) > \Phi(z)$: रखे गए खंड फेंके गए खंडों से
अधिक समृद्ध हैं, इसलिए हर कट-ऑफ़ चट्टान की तुलना में धातु का बड़ा भाग रखता है; यह अंतर
$s$ के साथ, अर्थात् निक्षेप की परिवर्तिता के साथ, बढ़ता है।
\end{solution}

\begin{solution}{exo:b2:ore-deposits:11}
(क) $D(R + 1)/(R + D) = fD$ से $R + 1 = fR + fD$, $R = (fD - 1)/(1 - f)$। (ख) $f = 0.5$: $R = D - 2$,
298 और \num{9998}; $f = 0.9$: $R = 9D - 10$, \num{2690} और \num{89990}। (ग) निकेल अपनी अधिकतम सीमा के
\qty{50}{\percent} तक निकेल--ताँबा अयस्कों के परास (100 से \num{2000}) के भीतर पहुँच जाता है, और 90 तक उसके ठीक ऊपर;
प्लैटिनम को $R$ के \num{10000} से \num{100000} तक के मानों की आवश्यकता होती है, जो प्लैटिनम-समूह तत्त्व रीफ़ों का परास है।
\end{solution}

\begin{solution}{exo:b2:ore-deposits:12}
(क) $\ln M = m + s u$ के साथ, क्वांटाइल से ऊपर के निक्षेपों के लिए $u > \Phi^{-1}(q)$; उनकी धातु
$\int_{\Phi^{-1}(q)}^{\infty} e^{m + su}\varphi(u)\,\mathrm{d}u = e^{m + s^2/2}[1 - \Phi(\Phi^{-1}(q) -
s)]$ है, और कुल $e^{m + s^2/2}$ से भाग देने पर $\Phi(s - \Phi^{-1}(q))$ मिलता है। (ख) $s = 0.702 \ln 10 =
1.616$; $\Phi(1.616 - 1.282) = \Phi(0.334) = 0.631$। (ग) $\Phi(1.616 - 2.326) = \Phi(-0.710) = 0.239$,
जबकि चार सबसे बड़े निक्षेपों के लिए 0.266: सबसे बड़े निक्षेप लघुगणकीय-प्रसामान्य नियम की अपेक्षा से
थोड़े बड़े हैं।
\end{solution}

\begin{solution}{pb:b2:ore-deposits:1}
\textbf{1.} \qty{330}{m}; $\sum l\,g = 138.4$, माध्य $138.4/330 = \qty{0.419}{\percent}$।
\textbf{2.} $3.24/8 = \qty{0.405}{\percent}$; हर अंतराल कोर की एक लंबाई का, अतः चट्टान के एक आयतन का,
प्रतिनिधित्व करता है: लंबे अंतराल को छोटे अंतराल से अधिक गिना जाना चाहिए।
\textbf{3.} $230/330 = 0.697$, $118.6/230 = \qty{0.516}{\percent}$ पर।
\textbf{4.} $4190/28 = 150$।
\textbf{5.} $8100/28 = 289$।
\textbf{6.} $0.38\,e^{0.55^2/2} = 0.38\,e^{0.1513} = \qty{0.442}{\percent}$; बोरहोल,
\qty{0.419}{\percent} पर, उससे \qty{6}{\percent} के भीतर है: प्रतिनिधिक।
\textbf{7.} $\pi \times 450^2 \times 700 = \qty{4.45e8}{m^3}$; $\times 2.7 = \num{1.20e9}$ t,
अर्थात् \qty{1200}{Mt}।
\textbf{8.} $z = \ln(0.38/0.3)/0.55 = 0.4298$; $\Phi(z) = 0.6663$।
\textbf{9.} $z = 1.167$, 0.4298, $-0.4990$; $\Phi = 0.8784$, 0.6663, 0.3089; $T = \qty{1050}{Mt}$,
\qty{800}{Mt}, \qty{371}{Mt}।
\textbf{10.} $\Phi(z + s) = \Phi(0.9798) = 0.8364$; $\bar g = 0.442 \times 0.8364/0.6663 =
\qty{0.555}{\percent}$।
\textbf{11.} $0.442 \times 0.9570/0.8784 = \qty{0.482}{\percent}$; $0.442 \times 0.5203/0.3089 =
\qty{0.744}{\percent}$।
\textbf{12.} \omterm{def:b2:ore-deposits:grade}{टनभार} $\times\, 370.7/799.6 = 0.464$; ग्रेड $\times\, 0.744/0.555 = 1.34$।
\textbf{13.} माध्यिका पर, \qty{0.38}{\percent} ($z = 0$): \qty{600}{Mt}, $0.442 \times
\Phi(0.55)/0.5 = 0.442 \times 0.7088/0.5 = \qty{0.627}{\percent}$ पर।
\textbf{14.} $1200 \times 0.00442 = \qty{5.30}{Mt}$।
\textbf{15.} $(\log_{10} 1200 - 2.373)/0.654 = (3.0792 - 2.373)/0.654 = 1.080$; $\Phi = 0.8599$: 86-वाँ
शतमक, एक बड़ा निक्षेप।
\textbf{16.} $(\log_{10} 0.442 + 0.364)/0.199 = (-0.3546 + 0.364)/0.199 = 0.047$; $\Phi = 0.5187$: एक
साधारण ग्रेड, 52-वाँ शतमक।
\textbf{17.} सहसंबंध के साथ 0.702, उसके बिना 0.684।
\textbf{18.} $\mu_M = 0.009$; $(\log_{10} 5.30 - 0.009)/0.702 = (0.7243 - 0.009)/0.702 = 1.019$; $\Phi =
0.8459$: 85-वाँ शतमक।
\textbf{19.} $422 \times (1 - 0.8459) = 65$।
\textbf{20.} $M = T_0\,e^{\mu + s^2/2}\,\Phi(z + s) = 1200 \times 0.00442 \times 0.9570 =
\qty{5.08}{Mt}$।
\textbf{21.} $1200 \times 0.00442 \times 0.5203 = \qty{2.76}{Mt}$।
\textbf{22.} ताँबे का $\Phi(z + s) = 0.8364$ भाग, \omterm{def:b2:ore-deposits:grade}{टनभार} के $\Phi(z) = 0.6663$ भाग में।
\textbf{23.} 0.3 से \qty{0.2}{\percent} तक \omterm{def:b2:ore-deposits:grade}{टनभार} $1054 - 800 = \qty{254}{Mt}$
(\qty{32}{\percent}) बढ़ता है और ताँबा $5.08 - 4.44 = \qty{0.64}{Mt}$ (\qty{14}{\percent}): जोड़े गए
खंडों का औसत $0.64/254 = \qty{0.25}{\percent}$ है, जो निक्षेप के सबसे निर्धन खंड हैं।
\textbf{24.} $1200 \times 0.00442 \times 0.8364 = 800 \times 0.00555$: पूर्वेक्षण-स्थल में \qty{0.3}{\percent}
कट-ऑफ़ से ऊपर \textbf{\qty[text-series-to-math=true]{4.44}{Mt}} ताँबा है।
\end{solution}
